% solution source: 2025_SKYMAP_sol.pdf pages 12-14 (OM03: Solution and Marking Scheme)
\textbf{(OM03.1)} Draw the celestial equator (the line passing through W and
Mintaka ($\delta$ Ori)). Measure the angle $\phi_X$ between the Celestial
Equator and the vertical at $270^\circ$ azimuth.
\[ \phi_X = 21^\circ\,\mathrm{N} \]
Full credit for $20^\circ$ to $22^\circ$; half credit at the limits
$19^\circ$ and $23^\circ$. \hfill(4 marks: value 1.0; drawing the equator
0.5; equator passing through Mintaka (point C) 1.0; indication of
hemisphere either through ``$+$'', ``N'' or ``North'' 1.5)

% FIGURE NEEDED: question grid with the celestial equator drawn through W
% and Mintaka and the angle phi_X = 21 deg marked at 270 deg azimuth
% (2025_SKYMAP_sol.pdf page 12)

\textbf{(OM03.2)} The star Mintaka (which lies very close to the Celestial
Equator, i.e.\ $\delta \approx 0^\circ$) has an approximate altitude at X at
22:00\,hrs on 21 March of $a_{X,\mathrm{21Mar},22:00} = 25^\circ$ from the
West horizon.

Therefore its hour angle at 22:00\,hrs local time on 21 March is given by
\[ H_{\mathrm{21Mar},22:00}
   = \cos^{-1}\!\left[\frac{\sin 25^\circ - \sin 0^\circ \sin 20^\circ}
                           {\cos 0^\circ \cos 20^\circ}\right]
   = 63.27^\circ \]
Therefore, the hour angle of Mintaka at 22:00\,hrs local time on 21 Jan 2026
(306 days later) will be
\[ H_{\mathrm{21Jan},22:00}
   = H_{\mathrm{21Mar},22:00} + 306 \times \frac{360^\circ}{365.2564}
   = 63.27^\circ + 301.60^\circ = 364.87^\circ \equiv 4.87^\circ \]
The hour angle of Mintaka at 18:00\,hrs local time on 21 Jan 2026 will be
\[ H_{\mathrm{21Jan},18:00}
   = H_{\mathrm{21Jan},22:00} - 4 \times \frac{360^\circ}{24}
   = 4.87^\circ - 60^\circ = -55.13^\circ \]
The hour angle of an object at a certain local time on a given day will be
the same at all locations.

Since the hour angle is negative, it indicates that Mintaka is to the east
of the meridian at 18:00\,hrs local time. Therefore, the point P represents
the cardinal point East.

The altitude of Mintaka at 18:00\,hrs local time on 21 Jan 2026 at the
location Y will be (since $\delta \approx 0$)
\[ a_{Y,\mathrm{21Jan},18:00}
   = \sin^{-1}\!\left[\cos H_{\mathrm{21Jan},18:00} \cos\phi_Y\right]
   = \sin^{-1}\!\left[\cos(-55.13^\circ)\cos(-40^\circ)\right]
   \approx 26^\circ \]
Now, the location Y lies in the Southern hemisphere. The Celestial Equator
lies towards the North (i.e.\ tilted towards the left of the cardinal point
E in the given projection) near the Eastern horizon in the Southern
hemisphere. Further, the latitude of Y is $40^\circ$\,S, which gives the
inclination of the Celestial Equator w.r.t.\ the vertical at $90^\circ$
azimuth. The position of Mintaka will be nearly on the Celestial Equator at
an altitude of $\approx 26^\circ$.

The orientation of Orion can be inferred by realizing that the shield of
Orion is ahead of the belt as the sky moves.

% FIGURE NEEDED: solution grid drawing of Orion at location Y with points
% A-F, the celestial equator inclined at 40 deg to the vertical, and P
% marked near the horizon (2025_SKYMAP_sol.pdf page 14)

Marking scheme (14 marks):
\begin{itemize}
\item ``P'' marked as ``E'' \hfill(2.0)
\item Angle of the line joining Mintaka (C) and cardinal point (P) w.r.t.\
  the vertical between $37^\circ$ and $43^\circ$ \hfill(2.0)
  (1.0 if between $35^\circ$--$37^\circ$ or $43^\circ$--$45^\circ$)
\item All 6 of the A--F points lie to the left of the vertical at $90^\circ$
  azimuth \hfill(2.0)
  (1.0 if only 5 points, 0.5 if only 4 points)
\item Altitude of Mintaka (point C) between $24^\circ$ and $28^\circ$
  \hfill(2.0)
  (1.0 if between $22^\circ$--$24^\circ$ or $28^\circ$--$30^\circ$)
\item Shield (point D) at lesser azimuth than belt (point C) \hfill(2.0)
\item Legs (points A and B) are south of the Celestial Equator \hfill(2.0)
\item Scale: distance B--F is between 1.8\,cm and 2.8\,cm \hfill(2.0)
  (1.0 if between 1.3\,cm--1.8\,cm or 2.8\,cm--3.3\,cm)
\end{itemize}
